\section{Progress Theorem}
\label{sec:progress}
$\minilang$ expressions are evaluated to values by
inputting them to the transition system defined
by the dynamic semantics given in Figure~\ref{fig:eval_rules}.
Each transition in the system \emph{reduces} an expression or
brings the expression closer to a value.
However, not every irreducible expression is a value such as
the following:

\texttt{+(num[5]; str[abc])} $\stepto \times$

An expression $e$ that is not a value, but for which there does
\emph{not} exists an $e'$ such that $e \stepto e'$ is said to be \emph{stuck}.
It should be the case that any stuck expression is ill-typed.
Moreover, well-typed expressions do not get stuck.
This property is expressed formally by the progress theorem.

\begin{thm}
\textbf{(Progress)}
If $e : \tau$, then either
\texttt{$e$ value} or there exists an expression $e'$ such
that $e \stepto e'$.
\end{thm}

Progress is proved by induction on the typing rules.
Again there a lot of cases, so for brevity
we show just one case to get an idea of
how to prove this theorem.

\subsection{Inductive Case: (\texttt{T.4}) -- Addition Case}

\[
\cinfer[T.4]{+($e_1$, $e_2$): num}
  {\code{$e_1$: num} & \code{$e_2$: num}}
\]
By the inductive hypothesis,
since \texttt{$e_1:$ num}, either \texttt{$e_1$ value}
or there exists an expression $e'_1$ s.t.
$e_1 \stepto e'_1$.

Suppose \texttt{$e_1$ value}. Then either
\texttt{$e_1=$  num[$n_1$]} or \texttt{$e_1=$  str[$s_1$]}.
Since \texttt{$e_1:$ num}, then it must be the case that
\texttt{$e_1=$  num[$n_1$]}.

Therefore, \texttt{$e_1=$  num[$n_1$]} or
there exists an expression $e'_1$ such that
$e_1 \stepto e'_1$.
Similarly, \texttt{$e_2=$  num[$n_2$]} or
$e_2 \stepto e'_2$ for some $e'_2$.

Suppose $e_1 \stepto e'_1$. Then:
\[
\cinfer[D.1]{+($e_1$; $e_2$) $\stepto$ +($e'_1$; $e_2$)}{e_1 \stepto e'_1}
\]

Suppose \texttt{$e_1=$  num[$n_1$]} and $e_2 \stepto e'_2$. Then:
\[
\cinfer[D.2]{+(num[$n_1$]; $e_2$) $\stepto$ +(num[$n_1$]; $e'_2$)}{e_2 \stepto e'_2}
\]

Suppose \texttt{$e_1=$  num[$n_1$]} and \texttt{$e_2=$  num[$n_2$]}. Then:
\[
\cinfer[D.3]{+(num[$n_1$]; num[$n_2$]) $\stepto$ num[$n_1 + n_2$]}{ }
\]

We have covered all of the nested cases for rule \texttt{T.4}. $\qed$
